\section{第四日第一讲（Akshay）}

\subsection{Plancherel 公式}

回忆我们有两个 TQFT：$A_G$ 与 $B_{\check G}$。取 $Y:=\mathbb F_q((t))$，我们期待一个 2-范畴的等价 $A_G(Y)\to B_{\check G}(Y)$。给定球 $G$-簇 $X$ 带来的额外数据，我们有相应的边界理论 $\Theta_{G,X}\in A_G(Y)$ 与 $\Theta_{\check G,\check X}\in B_{\check G}(Y)$。我们期待这两个 1-范畴也等价。

先分析我们想要的等价在去范畴化之后的版本。取 $F=\mathbb F_q((t))$、$O=\mathbb F_q[[t]]$，设 $X$ 是球 $G$-簇。在某些有限性假设下，我们有
\begin{equation}
  C_c(X(F))^{G(O)}\cong\mathbb C[\Rep\check G_X]
  \label{eq:11-1}
\end{equation}
其中 $\check G_X\subset\check G$ 是某个群。就本讲而言，我们设 $\check G_X=\check G$（这就是 BZSV 里所谓的 \textbf{strongly tempered} 条件）。上述同构不仅是 $H(G(F),G(O))\cong\mathbb C[\Rep\check G]$ 的模同构，它还保持内积。设 $f,g\in C_c(X(F))^{G(O)}$，$\hat f,\hat g$ 是它们在 (11.1) 下的像，则
\begin{equation}
  \langle f,g\rangle=\int_{\check G_c}\langle\hat f,\hat g\rangle\cdot\mathrm{q\text{-}char}(\check G\times G_{gr}^m,M^\vee)\,d\mu_H.
  \label{eq:11-2}
\end{equation}
记号说明如下：

\begin{itemize}[leftmargin=1.6em]
  \item $\check G_c$ 是复约化群 $\check G$ 的紧形式。
  \item $M^\vee$ 是某个带 $\check G\times G_{gr}^m$ 作用的空间，通常是 $T^*\check X$。
  \item $M^\vee$ 上的 q-特征标，指的是把 $\mathbb C[M^\vee]$ 看作 $\check G\times G_{gr}^m$ 的表示所得的 q-特征标。
  \item $\mu_H$ 是 Haar 测度。
\end{itemize}

在同构 (11.1) 下，示性函数 $\mathbf 1_{X(O)}$ 应当被送到平凡表示。

\begin{exmp}
我们已对下面两个情形算过 Plancherel 公式：
\begin{enumerate}[label=(\arabic*),leftmargin=1.8em]
  \item $G=\mathbb G_m$，$X=\mathbb A^1$。此时 $\check G=\mathbb G_m$，$M^\vee=T^*\mathbb A^1$。
  \item $G=\mathrm{PGL}_2$，$X=\mathrm{PGL}_2/\mathbb G_m$。此时 $\check G=\mathrm{SL}_2$，$M^\vee=T^*\mathbb A^2$。
\end{enumerate}
\end{exmp}

\subsection{内积作为 Hom 空间上的迹}

一个关键的观察是：(11.2) 的两端都可以实现为某些算子在某个 Hom 空间上的迹。设 $\hat f=\chi_V$（表示 $V$ 的特征标），$\hat g=\chi_W$。换句话说，$f=h_V*\mathbf 1_{X(O)}$，$g=h_W*\mathbf 1_{X(O)}$。为简单起见，我们把 (11.2) 的左端记作 LHS，右端记作 RHS。

注意若 $A,B$ 是 $\check G$ 的复表示，则 $\chi_A\chi_B$ 是 $\Hom(B,A)$ 的特征标，故
\begin{equation}
  \int_{\check G_c}\chi_A\chi_B\,d\mu_H=\dim\Hom(B,A)^{\check G}=\dim\Hom_{\check G}(B,A).
  \label{eq:11-4}
\end{equation}
于是
\begin{align}
  \mathrm{RHS}&=\int_{\check G_c}\chi_V\chi_W\Big(\sum_i q^{-\frac i2}\chi_{\mathbb C[M^\vee]_i}\Big)d\mu_H
  \label{eq:11-5}\\
  &=\sum_i q^{-\frac i2}\dim\Hom_{\check G}(W,V\otimes\mathbb C[M^\vee]_i)\,d\mu_H
  \label{eq:11-6}\\
  &=\operatorname{tr}\Big(q^{-\frac12},\Hom_{\check G}(W,V\otimes\mathbb C[M^\vee])\Big).
  \label{eq:11-7}
\end{align}
在这些等式中，我们把 $q^{-\frac12}$ 看作 $G_{gr}^m(\mathbb C)=\mathbb C^\times$ 中的元素。

现在把 $\mathbb C$-线性映射 $W\to V\otimes\mathbb C[M^\vee]$ 看作 $\mathbb C[M^\vee]$-线性映射 $W\otimes\mathbb C[M^\vee]\to V\otimes\mathbb C[M^\vee]$，可以得到一个更对称的表述。于是
\begin{equation}
  \mathrm{RHS}=\operatorname{tr}\Big(q^{-\frac12},\Hom_{\check G,\mathbb C[M^\vee]}(W\otimes\mathbb C[M^\vee],V\otimes\mathbb C[M^\vee])\Big).
  \label{eq:11-8}
\end{equation}
令 $\mathcal V=V\otimes\mathbb C[M^\vee]$。它是一个 $O(M^\vee)$-模，带 $\check G$ 作用，因此可以看作 $M^\vee/\check G$ 上的凝聚层。类似地定义 $\mathcal W$。于是
\begin{equation}
  \mathrm{RHS}=\operatorname{tr}\Big(q^{-\frac12},\Hom_{M^\vee/\check G}(\mathcal W,\mathcal V)\Big).
  \label{eq:11-9}
\end{equation}
这里的 $\Hom$ 是在 $M^\vee/\check G$ 上拟凝聚层（indcoherent？）的导出范畴中取的。

(11.2) 的左端可以用层-函数对应实现为迹。回忆一个一般事实：设 $Z$ 是 $\mathbb F_q$ 上的簇，$F,G$ 是 $Z_{\mathbb F_q}$ 上的混合 $\bar{\mathbb Q}_\ell$-层，则取 Frobenius 的迹得到 $Z(\mathbb F_q)$ 上的函数 $f,g$。我们有
\begin{equation}
  \operatorname{tr}\Big(\mathrm{Fr}^{-1}\,\Big|\,\Hom_{Z_{\mathbb F_q}}(F,G)\Big)=\sum_{z\in Z(\mathbb F_q)}f(z)g(z).
  \label{eq:11-10}
\end{equation}
注意这里的 $g$ 是 $DG$ 的迹（$G$ 的 Verdier 对偶）。例如取 $F=DG=\bar{\mathbb Q}_{\ell,Z}$（$Z$ 上的常值层），上式就化归 Lefschetz 不动点公式。

现在取 $F=\mathbb F_q((t))$、$O=\mathbb F_q[[t]]$。对很多定义在 $\mathbb F_q$ 上的「好」$X$，$X(F)/G(O)$ 是某个「合理的代数栈」$X_F/G_O$ 的 $\mathbb F_q$-点。例如当 $X=H$、$G=H\times H$ 时，$X(F)/G(O)$ 就是仿射 Grassmann 流形 $\mathrm{Gr}_G$ 的 $\mathbb F_q$-点，而它是一个 ind-scheme。

Hecke 代数 $H_G$ 在 $C_c(X(F))^{G(O)}$ 上的作用，可以升级为 Hecke 范畴在 $X_F/G_O$ 上的层的作用。这里的 Hecke 范畴是 $\mathrm{Gr}_G$ 上 $G(O)$-等变可构造混合 $\bar{\mathbb Q}_\ell$-层的导出范畴。记 $\mathrm{IC}_V$ 为几何 Satake 同构下与表示 $V$ 对应的反常层。令 $\delta_0:=j_*(\bar{\mathbb Q}_{\ell,X_O/G_O})$，其中 $j:X_O/G_O\to X_F/G_O$ 是闭嵌入（？）。则在层-函数对应下，LHS 中的 $f,g$ 分别是与 $\mathrm{IC}_W*\delta_0$ 与 $\mathrm{IC}_V*\delta_0$ 相联系的函数。

于是可以写
\begin{equation}
  \mathrm{LHS}=\operatorname{tr}\Big(\mathrm{Fr}^{-1}\,\Big|\,\Hom(\mathrm{IC}_W*\delta_0,\mathrm{IC}_V*\delta_0)\Big).
  \label{eq:11-11}
\end{equation}
这里的 $\Hom$ 在 $X_F/G_O$ 上可构造层的范畴中取值。

\subsection{局部猜想}

上一节的讨论强烈提示一个范畴等价（并不正确，尚需修正）：
\begin{equation}
  \mathrm{Shc}(X_F/G_O)\cong\QCoh(M^\vee/\check G).
  \label{eq:11-12}
\end{equation}
该等价应当满足：

\begin{itemize}[leftmargin=1.6em]
  \item $\delta_0$ 应被送到 $M^\vee/\check G$ 上的结构层 $\mathcal O$。
  \item 左端的 Frobenius 作用应当对应右端 $q^{\frac12}\in G_{gr}^m$ 的作用。
  \item 该等价与 Hecke 范畴的模结构相容（在导出 Satake 同构之下）。在 Abel 范畴的层面，$H_G$ 从左端以卷积作用，$\Rep(\check G)$ 从右端以张量积作用：对 $\check G$ 的任意表示 $V$，沿 $M^\vee/\check G\to*/\check G$ 拉回可得 $M^\vee/\check G$ 上的向量丛，$V$ 的作用就是与该向量丛作张量积。
\end{itemize}

该等价可以看作把 (11.1) 开头讨论的边界条件匹配起来。

上面这个等价按原样并不正确。在说修正之前，先做一点合理性检验。若该等价成立，则应当有同构
\begin{equation}
  \Hom_{\mathrm{Shc}(X_F/G_O)}(\delta_0,\delta_0)\cong\Hom_{M^\vee/\check G}(\mathcal O,\mathcal O).
  \label{eq:11-13}
\end{equation}
左端等于 $H^*(X_O/G_O)$。注意 $G_O/G$ 是 pro-unipotent 的，因此可缩，对上同调没有贡献。类似地 $X_O$ 收缩到 $X(\mathbb F_q)$。于是
\begin{equation}
  H^*(X_O/G_O)\cong H^*_G(X(\mathbb C)).
  \label{eq:11-14}
\end{equation}
右端则有
\begin{equation}
  \Hom_{\QCoh(M^\vee/\check G)}(\mathcal O,\mathcal O)\cong\mathbb C[M^\vee]^{\check G}.
  \label{eq:11-15}
\end{equation}
所以我们应当检验
\begin{equation}
  H^*_G(X(\mathbb C))=\mathbb C[M^\vee]^{\check G}.
  \label{eq:11-16}
\end{equation}

\begin{exmp}
考虑 $G=\mathbb G_m$、$X=\mathbb A^1$、$\check G=\mathbb G_m$、$M^\vee=T^*\mathbb A^1$ 的情形。此时 $H^*_G(X)=\mathbb C[\xi^2]$，其中 $\xi^2$ 位于同调次数 $2$。回忆 $\lambda\in\check G$ 作用在 $(x,y)\in T^*\mathbb A^1$ 上为 $(\lambda x,\lambda^{-1}y)$。我们有 $\mathbb C[M^\vee]^{\check G}=\mathbb C[xy]$。
\end{exmp}

\begin{exmp}
考虑 $G=\mathrm{PGL}_2$、$X=\mathrm{PGL}_2/\mathbb G_m$、$\check G=\mathrm{SL}_2$、$M^\vee=T^*\mathbb A^2$ 的情形。此时 $H^*_G(X)=H^*_{\mathbb G_m}(\mathrm{pt})=\mathbb C[\xi^2]$。另一方面，把 $M^\vee=T^*\mathbb A^2$ 实现出来，$\mathrm{SL}_2$ 以左乘作用，则有 $\mathbb C[M^\vee]^{\check G}=\mathbb C[\det]$。
\end{exmp}

两个例子中都验证了代数同构，但还得给右端配上正确的分次。在范畴层面有一个称为 \textbf{shearing} 的步骤，修正后的猜想如下：
\begin{equation}
  \mathrm{Shc}(X_F/G_O)\cong\IndCoh(M^\vee/\check G)^\natural.
  \label{eq:11-19}
\end{equation}
